% =====================================================================
% Thresholds, Risk and Settlement: a model of last-round draws in chess
% ---------------------------------------------------------------------
% Bibliography status (checked 2026-10-06 against RePEc records):
%  - Moul & Nye (2009), JEBO 70(1-2):10-21        : verified (RePEc)
%  - Krumer, Megidish & Sela (2023), MSS 122:50-57 : verified (RePEc)
%  - Vong (2017), GEB 102:562-567                  : verified (RePEc)
%  - Csato & Sziklai (2026), arXiv:2610.01997      : read in full
%  - Csato & Krumer, arXiv:2410.19333 (v3)         : abstract read
%  - Feddersen et al. (2023), Csato et al. (2024),
%    Csato & Gyimesi (2026, EJOR; football)        : from the reference
%    list of arXiv:2610.01997
% =====================================================================
\documentclass[11pt]{article}
\usepackage[margin=1in]{geometry}
\usepackage{amsmath,amssymb,amsthm,mathtools}
\usepackage{booktabs}
\usepackage[authoryear,round]{natbib}
\usepackage[pdftitle={Thresholds, Risk and Settlement: A Game-Theoretic Model of Last-Round Draws in Chess Tournaments},pdfauthor={Ambrose Birtwistle}]{hyperref}
\usepackage{microtype}

\newtheorem{proposition}{Proposition}
\newtheorem{corollary}{Corollary}
\newtheorem{lemma}{Lemma}
\theoremstyle{definition}
\newtheorem{assumption}{Assumption}
\newtheorem{remark}{Remark}

\newcommand{\Sset}{\mathsf{S}}   % settle
\newcommand{\Fset}{\mathsf{F}}   % fight
\newcommand{\EU}{\mathrm{EU}}

\title{Thresholds, Risk and Settlement:\\
A Game-Theoretic Model of Last-Round Draws in Chess Tournaments}
\author{Ambrose Birtwistle}
\date{October 2026 --- working draft}

\begin{document}
\maketitle

\begin{abstract}
In the final round of a chess tournament, players choose how much risk to
take, and the prizes, title norms and rating consequences they face are
step functions of the final score. I model a final-round pairing as a
two-player game in which each player either \emph{settles} for a draw or
\emph{fights} for a decisive result, and I show that the equilibrium depends
on a single statistic of each player's payoff: the share $\rho_i$ of the
payoff step that lies below the draw. A mutual-settlement region exists
if and only if $\rho_1+\rho_2>1$, so draw agreements require
aggregate concavity of payoffs around the current score and cannot arise with
linear payoffs and no effort cost. When fighting is costly and decisiveness
is complementary in the two players' efforts, settling and fighting are both
self-enforcing conventions: draw collusion needs no side payments, only a
focal point. The model yields testable implications for draw frequency,
game length and engine-measured risk as a function of the standings
before the last round.
\end{abstract}

\section{Introduction}

Chess differs from most tournament sports in one respect that matters for
incentives: a game has three outcomes, and either player can often steer
the game toward the middle one by simplifying the position. When prizes or
title norms are awarded at score thresholds, the middle outcome can be worth
nearly as much as a win to one player and nearly as little as a loss to
another. The consequences are familiar to organisers, who have introduced
rules such as bans on early draw offers and three-point scoring for wins, and
to the Qualification Commission of FIDE, which reports having investigated
tournaments where results were manipulated to obtain titles or norms and
lists short draws among its indicators \citep{fide_qc}.

The existing economic literature on strategic draws is, to my knowledge,
mostly empirical or simulation-based. The model here is complementary: it
characterises when threshold-type payoffs produce mutual settlement, when
they produce mutual risk-taking, and when both are equilibria.

\section{Related literature}

The closest empirical paper is \citet{moulnye}, who use post-war
tournament data to test whether Soviet players acted as a cartel by drawing
among themselves, and conclude that the pattern of draws is consistent
with collusion and inconsistent with competition. Their paper is a
statistical test for a specific group and period; it does not derive
how the shape of the prize function drives the choice between settling
and fighting.

Recent operations-research work studies chess tournament design.
\citet{csato_sziklai} simulate Swiss-system team tournaments to quantify the
trade-off between ranking accuracy and the probability of a stakeless last
round, and note that match-point scoring in the Chess Olympiad creates
situations where a draw on one board suffices for the team; they call for
empirical study of how extra points for winning affect behaviour.
\citet{csato_krumer} document that in Swiss tournaments with an odd number of
rounds, players with an extra game with the white pieces score more. These
papers concern design and fairness rather than individual strategic choice.

More generally, tournament theory has shown that contestants may rationally
shirk or even lose when future stages reward it \citep{vong2017,krumer2020}.
Stakeless last-round matches have been analysed in football
\citep{chater2021,csato2024,csato_gyimesi2026}, and \citet{feddersen2023} find
that betting markets anticipate effort adjustments at the end of the season.
I am not aware of a chess model with a continuous prize step function and a
draw outcome of the kind studied below, but I have not conducted a systematic
search, and this should be checked before any claim of novelty.

\section{The model}

\subsection{Setup}

Two players $i\in\{1,2\}$ are paired in the final round. Player $i$ has
current score $s_i$. Each chooses an action $a_i\in\{\Sset,\Fset\}$:
\emph{settle} (play to reduce variance and accept a draw when available) or
\emph{fight} (play for a decisive result). Fighting costs $c_i\ge 0$ in effort
or preparation. Write $a=(a_1,a_2)$.

\begin{assumption}[Decisiveness]\label{as:D}
The probability that the game is decisive is $D(a)\in(0,1)$, strictly
increasing in each player's action: $D(\Fset,a_2)>D(\Sset,a_2)$ and
$D(a_1,\Fset)>D(a_1,\Sset)$.
\end{assumption}

\begin{assumption}[Conditional win probability]\label{as:q}
Conditional on a decisive result, player~1 wins with probability
$q\in(0,1)$ and player~2 with probability $1-q$, independently of $a$.
Write $q_1=q$ and $q_2=1-q$.
\end{assumption}

The parameter $q$ is where the rating difference enters; for instance it can
be calibrated from a rating-based model of win, draw and loss probabilities
\citep{elo,milvang}. Holding $q$ fixed across actions is a simplification
discussed in Section~\ref{sec:limits}.

\begin{assumption}[Payoffs]\label{as:u}
Player $i$'s payoff from the tournament is $u_i(x)$, weakly increasing in the
final score $x$. A draw adds $\tfrac12$ and a win adds $1$ to $s_i$. Define
\[
u_{i0}=u_i(s_i),\quad u_{ih}=u_i(s_i+\tfrac12),\quad u_{i1}=u_i(s_i+1),
\]
\[
U_i^{+}=u_{i1}-u_{ih}\ge 0,\qquad U_i^{-}=u_{ih}-u_{i0}\ge 0 .
\]
Whenever $U_i^{+}+U_i^{-}>0$ define
\[
\rho_i=\frac{U_i^{-}}{U_i^{+}+U_i^{-}}\in[0,1].
\]
\end{assumption}

The statistic $\rho_i$ is the share of the payoff step lying
\emph{below} the draw: $\rho_i$ near $1$ means that a draw already secures
almost everything and a win adds little (a concave region), while $\rho_i$
near $0$ means that only a win matters (a convex region); $\rho_i=\tfrac12$
corresponds to linear payoffs.

\subsection{Expected payoff}

\begin{lemma}\label{lem:EU}
For $i\in\{1,2\}$,
\[
\EU_i(a)=u_{ih}+D(a)\,g_i-c_i\,\mathbf 1\{a_i=\Fset\},
\qquad g_i=q_iU_i^{+}-(1-q_i)U_i^{-}.
\]
Moreover $g_i>0\iff q_i>\rho_i$, and $g_i=0$ if $U_i^{+}+U_i^{-}=0$.
\end{lemma}

\begin{proof}
The player wins with probability $D q_i$, loses with probability
$D(1-q_i)$ and draws with probability $1-D$, so
$\EU_i+c_i\mathbf 1\{a_i=\Fset\}=u_{i0}D(1-q_i)+u_{ih}(1-D)+u_{i1}Dq_i
=u_{ih}+D\,[q_i(u_{i1}-u_{ih})-(1-q_i)(u_{ih}-u_{i0})]$.
Finally $g_i>0\iff q_i(U_i^{+}+U_i^{-})>U_i^{-}\iff q_i>\rho_i$.
\end{proof}

The number $g_i$ is the net value to player $i$ of adding variance: the gain
on the upside weighted by the chance of winning, minus the loss on the
downside weighted by the chance of losing.

\begin{lemma}[Draw probability by profile]\label{lem:draw}
The probability of a draw is $1-D(a)$. By Assumption~\ref{as:D},
$1-D(\Sset,\Sset)>1-D(\Sset,\Fset),\,1-D(\Fset,\Sset)>1-D(\Fset,\Fset)$.
\end{lemma}

\section{Results}

\subsection{Benchmark: no effort cost}

\begin{proposition}[Dominant strategies]\label{prop:dom}
Let $c_1=c_2=0$ and $g_i\neq 0$. Then fighting is strictly dominant for player
$i$ if $g_i>0$ and settling is strictly dominant if $g_i<0$. Hence the game has
a unique Nash equilibrium, in dominant strategies.
\end{proposition}

\begin{proof}
By Lemma~\ref{lem:EU}, for player~1,
$\EU_1(\Fset,a_2)-\EU_1(\Sset,a_2)=g_1\,[D(\Fset,a_2)-D(\Sset,a_2)]$, and
the bracket is strictly positive by Assumption~\ref{as:D} whatever $a_2$ is.
The sign is therefore that of $g_1$. The argument for player~2 is identical.
\end{proof}

\begin{proposition}[Settlement and fighting regions]\label{prop:regions}
Let $c_1=c_2=0$ and suppose $\rho_1,\rho_2$ are defined and $q\notin\{\rho_1,1-\rho_2\}$.
\begin{enumerate}
\item[(a)] $(\Sset,\Sset)$ is the equilibrium if and only if
$1-\rho_2<q<\rho_1$. This interval is nonempty if and only if
$\rho_1+\rho_2>1$, and its length is $\rho_1+\rho_2-1$.
\item[(b)] $(\Fset,\Fset)$ is the equilibrium if and only if
$\rho_1<q<1-\rho_2$. This interval is nonempty if and only if
$\rho_1+\rho_2<1$.
\item[(c)] $(\Fset,\Sset)$ is the equilibrium if and only if
$q>\max\{\rho_1,1-\rho_2\}$; $(\Sset,\Fset)$ if and only if
$q<\min\{\rho_1,1-\rho_2\}$.
\end{enumerate}
\end{proposition}

\begin{proof}
By Proposition~\ref{prop:dom} and Lemma~\ref{lem:EU}, player~1 fights iff
$q>\rho_1$ and player~2 fights iff $1-q>\rho_2$, i.e.\ $q<1-\rho_2$.
Intersecting the four sign combinations gives (a)--(c). The interval in (a) is
nonempty iff $1-\rho_2<\rho_1$, and (b) iff $\rho_1<1-\rho_2$.
\end{proof}

\begin{corollary}[Linear payoffs]\label{cor:linear}
If $U_i^{+}=U_i^{-}$ for both players (so $\rho_1=\rho_2=\tfrac12$), the
settlement region of Proposition~\ref{prop:regions} is empty: with no effort
cost, linear payoffs never generate a mutually settled draw. Mutual settlement
requires $\rho_1+\rho_2>1$, that is, payoffs that are on aggregate concave
around the current score.
\end{corollary}

\subsection{Threshold prizes}

Prizes, title norms and rating-performance cutoffs are step functions of
the final score. Let player $i$ receive a prize $V>0$ if and only if the final
score is at least $s_i+r_i$, plus a small linear component
$\varepsilon(x-s_i)$ with $\varepsilon>0$ for the rating or tie-break value of
each point:
\[
u_i(x)=V\,\mathbf 1\{x\ge s_i+r_i\}+\varepsilon\,(x-s_i).
\]

\begin{proposition}[Threshold requirements]\label{prop:threshold}
Let $r_i\in\{\le 0,\tfrac12,1,>1\}$ index how many additional points the player
needs.
\begin{enumerate}
\item[(i)] Draw suffices ($r_i=\tfrac12$): $U_i^{+}=\varepsilon/2$,
$U_i^{-}=V+\varepsilon/2$, and
$\rho_i=\dfrac{2V+\varepsilon}{2V+2\varepsilon}=1-\dfrac{\varepsilon}{2(V+\varepsilon)}$.
Player $i$ settles iff $q_i<\rho_i$.
\item[(ii)] Win needed ($r_i=1$): $U_i^{+}=V+\varepsilon/2$,
$U_i^{-}=\varepsilon/2$, and $\rho_i=\dfrac{\varepsilon}{2(V+\varepsilon)}$.
Player $i$ fights iff $q_i>\rho_i$.
\item[(iii)] Prize already secured or out of reach ($r_i\le0$ or $r_i>1$):
$U_i^{+}=U_i^{-}=\varepsilon/2$, $\rho_i=\tfrac12$, and player $i$ fights
iff $q_i>\tfrac12$.
\end{enumerate}
When $\varepsilon/V\to0$, in case (i) the player settles for every $q_i<1$,
and in case (ii) fights for every $q_i>0$.
\end{proposition}

\begin{proof}
Direct substitution into the definitions of $U_i^{\pm}$ and $\rho_i$, and
Lemma~\ref{lem:EU}. For (i), $u_{i0}=0$, $u_{ih}=V+\varepsilon/2$,
$u_{i1}=V+\varepsilon$. For (ii), $u_{i0}=0$, $u_{ih}=\varepsilon/2$,
$u_{i1}=V+\varepsilon$. For (iii), $V$ drops out of all differences.
\end{proof}

\begin{corollary}[Who settles]\label{cor:who}
If both players are in case~(i) (for instance, both share first place with
a draw), then $\rho_1+\rho_2>1$ and the settlement region is nearly all of
$(0,1)$, with length $\rho_1+\rho_2-1\to1$ as $\varepsilon/V\to0$. If one
player is in case~(i) and the other in case~(ii), then
$\rho_1+\rho_2=\tfrac{2V+\varepsilon}{2V+2\varepsilon}+\tfrac{\varepsilon}{2V+2\varepsilon}=1$:
the leader settles and the challenger fights for almost every $q$. If both
are in case~(iii) the only stake is the small linear component and the
conclusion depends on costs (Proposition~\ref{prop:cost}).
\end{corollary}

\begin{remark}[Scoring rules]\label{rem:scoring}
Under one point for a win and one half for a draw, the live requirements are
$r\in\{\tfrac12,1\}$, and a draw suffices in one of two cases. Under a
$3$--$1$--$0$ rule the live requirements are $r\in\{1,2,3\}$ points, and a draw
suffices in one of three. If live requirements were equally likely (an
illustrative assumption, not an empirical claim), the share of
draw-sufficient cases falls from $1/2$ to $1/3$, and the chance that both
players are simultaneously in case~(i), if independent, from $1/4$ to $1/9$.
This is consistent with the stated motivation of three-point rules, but the
rule also changes the values of $U_i^{\pm}$ themselves, so the comparison is
only heuristic.
\end{remark}

\subsection{Effort costs and tacit coordination}

\begin{assumption}[Complementarity]\label{as:comp}
Let $\Delta_1(a_2)=D(\Fset,a_2)-D(\Sset,a_2)$ and
$\Delta_2(a_1)=D(a_1,\Fset)-D(a_1,\Sset)$. Then
$\Delta_i(\Fset)\ge\Delta_i(\Sset)>0$: the effect of one's own fighting on
decisiveness is larger when the opponent also fights.
\end{assumption}

\begin{proposition}[Costly fighting]\label{prop:cost}
Let $c_i>0$ and, for each $i$ with $U_i^{+}+U_i^{-}>0$, define
\[
\rho_i^{\Sset}=\rho_i+\frac{c_i}{\Delta_i(\Sset)(U_i^{+}+U_i^{-})},\qquad
\rho_i^{\Fset}=\rho_i+\frac{c_i}{\Delta_i(\Fset)(U_i^{+}+U_i^{-})},
\]
so that $\rho_i\le\rho_i^{\Fset}\le\rho_i^{\Sset}$. Under
Assumption~\ref{as:comp}, and generically:
\begin{enumerate}
\item[(a)] Fighting is strictly dominant for $i$ iff $q_i>\rho_i^{\Sset}$,
and settling is strictly dominant iff $q_i<\rho_i^{\Fset}$.
\item[(b)] If $\rho_i^{\Fset}<q_i<\rho_i^{\Sset}$, player $i$ plays a
\emph{conditional} strategy: fight if the opponent fights, settle otherwise.
\item[(c)] If both players are conditional, the pure-strategy equilibria are
exactly $(\Sset,\Sset)$ and $(\Fset,\Fset)$.
\end{enumerate}
In general, the pure equilibria are the profiles in which every dominant player
plays her dominant action and every conditional player copies the opponent.
\end{proposition}

\begin{proof}
By Lemma~\ref{lem:EU}, the gain to player $i$ from fighting is
$G_i(a_j)=g_i\Delta_i(a_j)-c_i$. If $g_i\le0$ then $G_i<0$ for both $a_j$,
so settling is dominant (and $q_i\le\rho_i<\rho_i^{\Fset}$). If $g_i>0$,
$G_i(\Fset)\ge G_i(\Sset)$ by Assumption~\ref{as:comp}. Hence fighting is
dominant iff $G_i(\Sset)>0$, settling is dominant iff $G_i(\Fset)<0$, and
otherwise $G_i(\Sset)\le 0\le G_i(\Fset)$ (conditional). Solving
$g_i\Delta_i(a_j)>c_i$ for $q_i$ using $g_i=q_i(U_i^{+}+U_i^{-})-U_i^{-}$
gives the thresholds $\rho_i^{a_j}$. Part (c) follows because each player's
best response is then to match the other's action.
\end{proof}

\begin{corollary}[Collusion as coordination]\label{cor:coord}
When both players are conditional, settling and fighting are both
self-enforcing. Mutual settlement does not require an agreement enforceable by
side payments: it is a convention, and what selects it is whatever makes the
draw a focal point for the two players (a shared team, national federation,
club, friendship or reputation). Effort costs and complementarity widen the
region in which settlement is an equilibrium relative to
Proposition~\ref{prop:regions}.
\end{corollary}

\section{Empirical implications}

The model's observable object is the standings before the last round. They
determine each player's requirement $r_i$ (Proposition~\ref{prop:threshold}),
hence $\rho_i$ and the predicted regime.

\begin{enumerate}
\item[\textbf{H1}] The probability of a draw is higher when
$\rho_1+\rho_2>1$ than when $\rho_1+\rho_2<1$ (Propositions~\ref{prop:regions}
and Lemma~\ref{lem:draw}).
\item[\textbf{H2}] Holding ratings fixed, the draw probability changes
sharply at the score margin at which a draw starts to suffice
(case~(i)), because $\rho_i$ jumps between cases~(ii) and~(i).
\item[\textbf{H3}] Within the settlement region, draws are more frequent
between players who share an affiliation (Corollary~\ref{cor:coord}), the
channel studied empirically by \citet{moulnye} for a specific group.
\item[\textbf{H4}] Risk measured by an engine (for example the volatility of
the engine evaluation along the game, or the loss in centipawns) is lower under
mutual settlement and higher under mutual fighting.
\end{enumerate}

The data requirement is tournament-level results with final-round pairings,
scores, ratings and, ideally, prize or norm structures. Results sources such
as chess-results.com are used in the Swiss-system literature
\citep{csato_sziklai}; their terms of use should be checked before collection.
Draw offers are generally not recorded in over-the-board data, so game length
is the natural observable proxy.

\section{Extension: draw offers as signals (sketch)}

The static model treats the game's decisiveness as an exogenous function of
actions. A natural next step lets player~1 hold private information about the
objective evaluation of the position, $\theta\in\{\text{good},\text{bad}\}$,
and lets her offer or withhold a draw. The opponent updates on the offer and
accepts or refuses, and the payoff step $\rho_i$ determines which types
are willing to offer. This has the structure of a signalling game, as in a
prior paper of mine on programmatic signalling, but I have \emph{not} derived
or verified any result for it here, and it is left as a direction for a
second paper.

\section{Limitations}\label{sec:limits}

The model is static with binary actions and takes $q$ as independent of
behaviour. If instead rating-implied expected scores were held fixed whatever
the players do, a decisive game would shift $q$ with $D$, and the
region boundaries would move; the qualitative role of $\rho_i$ would remain
but Propositions~\ref{prop:dom}--\ref{prop:regions} would need to be restated.
Payoffs are reduced-form functions of the score and ignore
tie-break systems, team incentives and the possibility that draws affect
others in the standings. The model also covers a single pairing, so it ignores
the interaction between simultaneous boards (relevant to prizes decided by
several games).

\section*{Appendix: numerical check}

I verified Propositions~\ref{prop:dom}, \ref{prop:regions}--\ref{prop:threshold}
and~\ref{prop:cost} by brute-force search over pure-strategy equilibria for
randomly drawn $D$, $q$, payoff steps and costs (several hundred thousand
draws each, with random weakly increasing $u_i$ including flat steps), and
found no counterexample. This is a sanity check on the algebra, not a
substitute for the proofs above.

\begin{thebibliography}{99}
\bibitem[Chater et~al.(2021)]{chater2021}
Chater, M., Arrondel, L., Gayant, J.-P., and Laslier, J.-F. (2021).
Fixing match-fixing: Optimal schedules to promote competitiveness.
\emph{European Journal of Operational Research}, 294(2):673--683.

\bibitem[Csat\'o et~al.(2024)]{csato2024}
Csat\'o, L., Molontay, R., and Pint\'er, J. (2024).
Tournament schedules and incentives in a double round-robin tournament with
four teams. \emph{International Transactions in Operational Research},
31(3):1486--1514.

\bibitem[Csat\'o and Gyimesi(2026)]{csato_gyimesi2026}
Csat\'o, L. and Gyimesi, A. (2026).
Increasing competitiveness by imbalanced groups: The example of the 48-team
FIFA World Cup. \emph{European Journal of Operational Research},
332(3):868--879.

\bibitem[Csat\'o and Krumer(2026)]{csato_krumer}
Csat\'o, L. and Krumer, A. (2026).
Swiss-system chess tournaments and unfairness. arXiv:2410.19333 (v3).

\bibitem[Csat\'o and Sziklai(2026)]{csato_sziklai}
Csat\'o, L. and Sziklai, B.~R. (2026).
The efficacy-competitiveness frontier in the Swiss system. arXiv:2610.01997.

\bibitem[Elo(1978)]{elo}
Elo, A.~E. (1978). \emph{The Rating of Chessplayers, Past and Present}.
Arco, New York.

\bibitem[Feddersen et~al.(2023)]{feddersen2023}
Feddersen, A., Humphreys, B.~R., and Soebbing, B.~P. (2023).
Contest incentives, team effort, and betting market outcomes in European
football. \emph{European Sport Management Quarterly}, 23(3):605--621.

\bibitem[FIDE Qualification Commission(2024)]{fide_qc}
FIDE Qualification Commission (2024). Policy on Ratings and Titles.
\url{https://www.fide.com/2024/01/31/}.

\bibitem[Krumer et~al.(2023)]{krumer2020}
Krumer, A., Megidish, R. and Sela, A. (2023). Strategic manipulations in
round-robin tournaments. \emph{Mathematical Social Sciences}, 122:50--57.

\bibitem[Milvang(2015)]{milvang}
Milvang, O. (2015). Probability for the outcome of a chess game based on
rating. Manuscript.

\bibitem[Moul and Nye(2009)]{moulnye}
Moul, C.~C. and Nye, J.~V.~C. (2009). Did the Soviets collude? A statistical
analysis of championship chess 1940--1978. \emph{Journal of Economic Behavior
\& Organization}, 70(1--2):10--21.

\bibitem[Vong(2017)]{vong2017}
Vong, A.~I.~K. (2017). Strategic manipulation in tournament games.
\emph{Games and Economic Behavior}, 102:562--567.

\end{thebibliography}
\end{document}
